\chapter{Цель работы}

Подключить PostgreSQL к персональному ассистенту, проверить сохранность
заметок при пересоздании контейнеров, подключить файл конфигурации с хоста
и проверить изоляцию сети Docker.

\section{Персональные параметры}

\begin{table}[htbp]
\centering
\caption{Параметры запуска}
\begin{tabularx}{\textwidth}{@{}Xl@{}}
\hline
Параметр & Значение \\
\hline
Номер ИСУ & 467857 \\
Порт PostgreSQL на хосте & $5430+57=5487$ \\
Порт ассистента & $8000+57=8057$ \\
Префикс ресурсов & \texttt{fedchenko} \\
Учебный ключ приложения & \texttt{secret57fedchenko} \\
Имя базы данных & \texttt{fedchenko57} \\
\hline
\end{tabularx}
\end{table}

\chapter{Ход работы}

Команды выполняются из каталога \texttt{solutions/assist-py}.
Сценарий \texttt{lab3.sh} содержит отдельную команду для каждого этапа.

\newcommand{\LabEvidence}[2]{%
  \begin{figure}[htbp]
    \centering
    \IfFileExists{#1.png}{%
      \includegraphics[width=\linewidth,height=0.65\textheight,keepaspectratio]{#1.png}%
    }{%
      \fbox{\parbox{0.85\linewidth}{Скриншот ожидается: \texttt{#1.png}.}}%
    }
    \caption{#2}
  \end{figure}
}

\section{Именованный том и PostgreSQL}

Команда \texttt{bash lab3.sh 1} создаёт том \texttt{fedchenko-pgdata} и
подключает его к каталогу данных PostgreSQL. На хосте база доступна через
порт 5487, который перенаправляется во внутренний порт 5432. Ассистент
обращается к базе по адресу \texttt{host.docker.internal:5487}. Проверочный
вывод включает главную страницу, состояние БД, создание заметки и чтение
списка заметок.

\LabEvidence{step-01-curl}{Четыре запроса к ассистенту}
\clearpage

\section{Персистентность при пересоздании}

Команда \texttt{bash lab3.sh 2} останавливает и удаляет контейнер ассистента,
запускает новый контейнер с тем же адресом базы и проверяет, что
идентификатор контейнера изменился. Повторный запрос к \texttt{/notes}
проверяет наличие первой заметки.

\LabEvidence{step-02-persistence}{Проверка заметок после пересоздания ассистента}
\clearpage

\section{Привязка каталога конфигурации}

Команда \texttt{bash lab3.sh 3} подключает каталог \texttt{config} с хоста
в \texttt{/app/config} в режиме только для чтения. Через \texttt{docker exec}
выводится файл \texttt{app.conf}, где поле владельца содержит
\texttt{fedchenko}.

\LabEvidence{step-03-config}{Файл конфигурации внутри контейнера}
\clearpage

\section{Пользовательская сеть}

Команда \texttt{bash lab3.sh 4} пересоздаёт контейнеры в сети
\texttt{fedchenko-backend} и повторно подключает тот же том. Приложение
использует имя \texttt{fedchenko-postgres} и внутренний порт 5432.
Порт PostgreSQL на хосте не публикуется. Проверка сравнивает доступность
имени контейнера из стандартной и пользовательской сетей Docker.

\LabEvidence{step-04-ping-outside}{Проверка имени PostgreSQL вне пользовательской сети}
\LabEvidence{step-04-ping-inside}{Доступ к PostgreSQL из пользовательской сети}
\clearpage

\section{Финальное состояние}

Команда \texttt{bash lab3.sh 5} выводит время, пользователя и хост,
список запущенных контейнеров, именованный том и пользовательскую сеть.

\LabEvidence{step-05-infrastructure}{Контейнеры, том и сеть}
\clearpage

\chapter{Вывод}

Хранение заметок вынесено в PostgreSQL. Именованный том сохраняет данные
при пересоздании контейнера базы. Привязка каталога позволяет обновить
файл конфигурации на хосте без пересборки образа. В пользовательской сети
ассистент обращается к PostgreSQL по DNS-имени контейнера.

\section*{Репозиторий}
\href{\RepositoryURL}{mlproger/maxim-k8s, ветка lab-3}
